% Practical Kaggle workflow. A proposed training path, not an executed training recipe.
\section{Training with Kaggle resources}
\sectionframe{04}{Planning our first\\training experiment}

\begin{frame}{A training workspace your team can open in Kaggle}
\begin{columns}[T,onlytextwidth]
\column{.47\textwidth}\eyebrow[teal]{PREPARE THE NOTEBOOK}
\begin{enumerate}\setlength\itemsep{3mm}\small
\item Keep the working starter for evaluation.
\item Create a training notebook for this competition.
\item Choose L4 $\times$ 4 and Internet off.
\end{enumerate}
\column{.48\textwidth}\eyebrow{ADD INPUT BEFORE STARTING}
\softpanel{orange}{Bring three ingredients}{Checked training examples\par Compatible Gemma weights\par Offline training packages}
\end{columns}
\vspace{4mm}
\takeaway{First check model and package compatibility, then complete one small training update.}
% Sources: official Upgraded Accelerators and Model Selection/Budget rules, archived via Kaggle API; official starter; PEFT quantization documentation.
\note{Kaggle provides competition-attached L4 x4 notebooks with Internet disabled. Use Add Input to attach the competition and required datasets/models before starting the GPU session. Keep the tested evaluation notebook intact and build a separate training environment. The official evaluation wheelhouse is not a verified SFT environment. The submitted agent must still use gemma-4-31b-it-qat-w4a16-ct. Do not assume that this compressed inference checkpoint can be trained directly; first verify the training loader, Gemma architecture, adapter target layers, and export compatibility. Four GPUs require an appropriate distributed training setup; seeing four devices does not automatically pool their memory. Preparing files can be done before GPU time is spent.}
\end{frame}

\begin{frame}{What does one training example look like?}
\tagline{Teach the next useful action from a checked repair trace.}
\begin{columns}[T,onlytextwidth]
\column{.48\textwidth}\eyebrow[teal]{CONTEXT}
\softpanel{teal}{Issue + earlier tool output}{
URL highlighting breaks at @.\par\vspace{2mm}
File search returned:\par
\pathcode{rich/highlighter.py}}
\column{.48\textwidth}\eyebrow{NEXT ACTION}
\softpanel{orange}{The training target}{
Call \pathcode{read_file}\par\vspace{2mm}
with \pathcode{filepath} set to\par
\pathcode{rich/highlighter.py}}
\end{columns}
\vspace{2mm}
\small SFT uses tool results as context and trains on checked assistant actions.\par\vspace{2mm}
\small With LoRA, training updates a small adapter while the base model stays fixed.
% Sources: official Data description; https://huggingface.co/docs/trl/en/sft_trainer (dataset formats and tool calling).
\note{The public package supplies issues, code snapshots, reference patches, and tests, but not a ready-made corpus of complete successful agent trajectories. Collect traces from the training split and check their patches with valid controls. A trace records the messages, tool requests, observations, and final changes. TRL supports conversational messages with tool calls and tool-role responses, plus a tools schema column. Verify the Gemma chat template and loss masking: labels should target the intended assistant outputs, not train the model to invent tool results. Reference patches can supervise training examples; held-out answers must never become inference inputs. A direct issue-to-patch dataset is a simpler experiment, but it does not teach the whole tool loop by itself.}
\end{frame}

\begin{frame}{Planned: our first small training run}
\tagline{TRL handles training; PEFT trains and exports the LoRA adapter.}

\begin{tikzpicture}
\path[use as bounding box] (0,-2.5) rectangle (13.35,.35);
\foreach \x/\c in {0/teal,3.5/orange,7/orange,10.5/teal}{\fill[\c!8,rounded corners=2mm] (\x,-2.3) rectangle +(2.95,2.5);}
\foreach \x in {0,3.5,7}{\draw[flow link] (\x+3,-1)--(\x+3.45,-1);}
\node[anchor=north west,text width=2.5cm,align=left] at (.2,0){\textbf{One update}\\[3mm]\small Small batch\\Check memory};
\node[anchor=north west,text width=2.5cm,align=left] at (3.7,0){\textbf{A tiny pilot}\\[3mm]\small 10--20 steps\\Watch the loss};
\node[anchor=north west,text width=2.5cm,align=left] at (7.2,0){\textbf{Save + reload}\\[3mm]\small Export weights\\Fresh session};
\node[anchor=north west,text width=2.5cm,align=left] at (10.7,0){\textbf{Compare}\\[3mm]\small Fresh tasks\\Same time limit};
\end{tikzpicture}

\vspace{2mm}
\small\textbf{Loss} measures fit to training examples. Test repair success on fresh tasks.\par
\vspace{3mm}
\begin{columns}[T,onlytextwidth]
\column{.48\textwidth}\eyebrow{SAVE}
\small Write to \pathcode{/kaggle/working/}.\par
Save Version with outputs before stopping.
\column{.48\textwidth}\eyebrow[teal]{RE-EVALUATE}
\small Attach the adapter to the official starter.\par Compare the same held-out tasks.
\end{columns}
% Sources: official starter for Kaggle paths and packaging; official model/adapter rules; https://huggingface.co/docs/trl/en/sft_trainer and https://huggingface.co/docs/peft/en/developer_guides/quantization.
\note{First verify one forward/backward update with a small batch and short context; inspect gradients and GPU memory. A proposed initial pilot is only 10--20 optimizer steps, not a full training run. Save both recovery state and an exported PEFT adapter, then reload the export in a fresh evaluation session. Store inputs under /kaggle/input (read-only) and outputs under /kaggle/working. Save a notebook version with its outputs or download/persist the artifact before ending the session; the live working directory is not a permanent team backup. A valid adapter directory needs adapter_config.json and adapter_model.safetensors. Reattach that artifact to the official evaluator, verify it really loads on the required quantized base, and compare adapter off/on using the same held-out tasks and budget. Only then package submission.zip with agent.yaml at its root; the unpacked package must remain below 3 GiB. Check GPU quota before each run and stop the GPU session after saving; L4 quota currently costs twice the older reference rate. These steps describe planned work, not a claim that a training notebook has already passed.}
\end{frame}
